% /Users/pikachu/books/payments-book/src/parts/part6/ch30-primary-sources.tex
\chapter{Как читать первоисточники}
\label{ch:primary-sources}

\begin{kaobox}[title=Что вы узнаете из этой главы]
  Где~искать первоисточники по~карточным платежам, EMV,
  PCI~DSS, правилам \scheme{Visa} и~\scheme{Mastercard},
  документам ЦБ~РФ и~НСПК; какие из~них открыты, какие
  доступны только участникам, и~как~работать с~этим~gap;
  как~отслеживать изменения в~стандартах и~правилах,
  чтобы не~узнавать о~них из~штрафов.
\end{kaobox}

Читать спецификацию~--- не~то~же самое, что читать учебник или
маркетинговый white paper. Учебник объясняет мотивацию и~ведёт
читателя за~руку; спецификация фиксирует обязательства: что разрешено,
что запрещено и~какой именно текст имеет нормативную силу. Поэтому
в~этой главе нас~интересуют не~только адреса документов, но~и~метод
чтения: в~каком порядке в~них входить, как отличать normative text
от~explanatory text и~как жить с~тем, что часть отрасли спрятана
за~NDA.

Предыдущие двадцать восемь глав этой книги описывают
устройство платёжного стека~--- протоколы, стандарты, правила,
архитектуру. Но~платёжная индустрия не~стоит на~месте:
карточные сети обновляют правила раз в~полгода, PCI~SSC
выпускает новые версии стандартов, ЦБ публикует положения
и~указания, EMVCo пересматривает спецификации. Книга фиксирует
состояние на~момент написания; чтобы оставаться актуальным,
нужно уметь работать с~первоисточниками напрямую.

Эта глава~--- навигационная карта: не~пересказ документов,
а~указатель, где~что~искать, как~читать и~как~не~пропустить
важное изменение.


\section{EMVCo}
\label{sec:ps-emvco}

EMVCo~--- консорциум, управляющий семейством стандартов EMV:
чиповые карты, бесконтактные платежи, токенизация, 3-D Secure,
QR-коды. Основан шестью платёжными системами (Visa, Mastercard,
JCB, American Express, China UnionPay, Discover), и~его~спецификации
де-факто определяют, как~работает карточный чип на~всей
планете~\sidecite{emvco-specs}.

\subsection*{Ключевые документы}

\textbf{EMV Chip Specifications}~--- четыре книги, описывающие
полный цикл EMV-транзакции~(см.~главу~\ref{ch:emv}):
Book~1~(Application Independent ICC to Terminal Interface
Requirements), Book~2~(Security and Key Management),
Book~3~(Application Specification), Book~4~(Cardholder,
Attendant, and Acquirer Interface Requirements). Доступны
бесплатно на~сайте EMVCo после регистрации.

\textbf{EMV Contactless Specifications}~--- отдельный набор
книг для~бесконтактных платежей: Book~A~(Architecture and
General Requirements), Book~B~(Entry Point), Book~C
(Kernel Specifications, по~одной на~каждую платёжную
систему)~(см.~главу~\ref{ch:contactless}).

\textbf{EMV Payment Tokenisation Specification}~--- стандарт
сетевой токенизации, определяющий архитектуру TSP, формат
DPAN, Assurance Level и~domain
restriction~(см.~главу~\ref{ch:tokenization}).

\textbf{EMV 3-D Secure Protocol and Core Functions
  Specification}~--- протокол 3DS~v2, включающий описание
трёх доменов, AReq/ARes, challenge flow, frictionless flow
и~decoupled authentication~(см.~главу~\ref{ch:3ds}).

\subsection*{Как читать}

Спецификации EMVCo объёмны (Book~3 одна занимает более 200~страниц)
и~написаны в~стиле инженерного стандарта: точно, но~без~пояснений
мотивации. Несколько рекомендаций для~первого чтения:

Сразу полезно различать два слоя текста. \emph{Normative text}
отвечает на~вопрос, что карта, терминал или участник \emph{обязаны}
делать; \emph{explanatory text} помогает понять замысел, но~не~заменяет
само требование. При~споре между \enquote{кажется, авторы имели в~виду}
и~буквальной формулировкой выигрывает нормативный текст.

\begin{itemize}
  \item Начинайте не~с~Book~1, а~с~Book~4: она~описывает
        пользовательские сценарии (что~видит держатель карты
        и~кассир) и~даёт контекст для~технических деталей
        в~остальных книгах.
  \item В~Book~3 сосредоточьтесь на~главах о~Transaction Flow
        и~Data Elements: именно они~описывают, что~происходит
        при~каждой транзакции.
  \item Tag-Length-Value (TLV)~--- базовый формат данных EMV;
        без~понимания TLV чтение Book~2 и~Book~3 будет
        затруднено. Структура TLV описана в~Book~3,
        Annex~B.
  \item EMVCo публикует бюллетени (EMV Specification Bulletins),
        которые вносят точечные изменения в~спецификации;
        читать бюллетени~--- единственный способ узнать
        об~изменениях между мажорными версиями.
\end{itemize}

\begin{practice}
  Спецификации EMVCo доступны бесплатно, но~требуют регистрации
  на~сайте \domain{emvco.com}. Регистрация открытая, без~проверки
  принадлежности к~финансовой организации. Скачивайте все~четыре
  книги и~contactless-спецификации сразу~--- они~ссылаются друг
  на~друга, и~читать одну без~остальных затруднительно.
\end{practice}


\section{PCI SSC}
\label{sec:ps-pci}

PCI Security Standards Council (PCI~SSC)~--- организация,
разрабатывающая и~поддерживающая стандарты безопасности
платёжных данных. Основана теми~же пятью сетями
(Visa, Mastercard, American Express, Discover, JCB)
в~2006~году~\sidecite{pci-dss-4}.

\subsection*{Ключевые документы}

\textbf{PCI~DSS}~--- основной стандарт, 12~требований
к~защите данных держателей
карт~(см.~главу~\ref{ch:pci-dss}). Текущая версия~---
4.0.1. Документ объёмный (около~360~страниц), но~хорошо
структурирован: каждое требование содержит описание,
цель, процедуру тестирования и~руководство.

\textbf{SAQ (Self-Assessment Questionnaire)}~--- опросники
самооценки, по~одному на~каждый тип интеграции (SAQ~A,
SAQ~A-EP, SAQ~B, SAQ~C, SAQ~D и~другие). SAQ~--- сокращённая
версия PCI~DSS, применимая к~мерчанту в~зависимости
от~того, как~он~обрабатывает карточные
данные~\sidecite{pci-saq}.

\textbf{PCI PIN Security Requirements}~--- стандарт защиты
PIN, применимый к~процессинговым центрам, эквайерам
и~всем, кто~обрабатывает PIN
block~(см.~главу~\ref{ch:processing-center}).

\textbf{PCI P2PE (Point-to-Point Encryption)}~--- стандарт
шифрования данных от~терминала до~точки
дешифрования~\sidecite{pci-p2pe}.

\subsection*{Как читать}

PCI~DSS~--- один из~немногих стандартов в~индустрии,
который можно читать последовательно: он~идёт от~общих
принципов (требование~1~--- сетевая безопасность)
к~конкретным (требование~12~--- политики безопасности),
и~каждое требование раскладывается на~под-требования
с~тестовыми процедурами.

\begin{itemize}
  \item Начните с~раздела \enquote{Scope of PCI DSS
          Requirements}: правильное определение CDE~--- половина
        работы по~compliance.
  \item Для~каждого требования читайте колонку
        \enquote{Guidance}~--- она~объясняет мотивацию,
        которой нет в~самом~тексте требования.
  \item PCI~SSC публикует \enquote{Information Supplements}~---
        дополнительные руководства по~конкретным темам
        (токенизация, облако, мобильные платежи). Они~бесплатны
        и~часто полезнее, чем~основной стандарт, для~решения
        конкретной задачи.
\end{itemize}

Все~документы PCI~SSC доступны бесплатно. Базовая
точка входа~--- \domain{pcisecuritystandards.org};
обычно достаточно простой регистрации.


\section{\scheme{Visa}}
\label{sec:ps-visa}

Документация \scheme{Visa}~--- одна из~самых обширных
в~индустрии и~одновременно одна из~самых труднодоступных:
часть документов публична, часть доступна только участникам
сети (member banks), а~часть~--- только по~NDA.

\subsection*{Публичные документы}

\textbf{Visa Core Rules and Visa Product and Service Rules}~---
основной свод правил, определяющий обязанности эмитентов,
эквайеров и~мерчантов. Публикуется дважды в~год (апрель
и~октябрь). Доступен бесплатно на~\domain{visa.com}
в~разделе~Visa
Rules~\sidecite{visa-core-rules}.

\textbf{Visa Dispute Management Guidelines}~---
руководство по~диспутам. Здесь собраны reason codes,
сроки, compelling evidence и~пороги
VDMP/VFMP~\sidecite{visa-dispute-rules}.

\textbf{Visa Supplemental Requirements}~--- дополнения
к~Core Rules для~конкретных регионов и~продуктов.

\subsection*{Документы с~ограниченным доступом}

\textbf{Visa Technical Acceptance Device Guide (TADG)}~---
техническое руководство по~приёму карт Visa на~терминалах;
содержит требования к~EMV-конфигурации, бесконтактным
платежам, обработке PIN. Доступен только для~участников
сети~\sidecite{visa-tadg}.

\textbf{Visa BASE~I / BASE~II Technical Specifications}~---
форматы сообщений для~авторизации и~клиринга. Доступны
только для~процессоров с~прямым подключением к~VisaNet.

\textbf{VTS (Visa Token Service) Technical Documentation}~---
спецификации токенизации. Доступны для~Token
Requestors~\sidecite{visa-vts}.

\subsection*{Как работать с~gap}

Если~вы~не~являетесь участником сети \scheme{Visa}
и~не~имеете доступа к~TADG или~BASE~I/II спецификациям,
Core Rules~--- ваш~основной источник. Они~не~содержат
технических деталей протокола, но~описывают все~бизнес-правила:
что~обязан делать эквайер, что~обязан делать эмитент,
какие транзакции допустимы, какие~--- нет.

Для~технических деталей, недоступных в~публичных документах,
работают два~подхода: (1)~запросить документацию через
вашего банка-эквайера или~процессора, который имеет
членство в~сети; (2)~использовать Visa Developer Portal
(\domain{developer.visa.com}), где~доступны API-спецификации,
sandbox-среды и~руководства по~интеграции для~конкретных
продуктов (Visa Direct, VTS, Visa Advanced Authorization).


\section{\scheme{Mastercard}}
\label{sec:ps-mastercard}

Документация \scheme{Mastercard} организована аналогично
\scheme{Visa}: публичные правила, ограниченные технические
спецификации и~developer-портал.

\subsection*{Публичные документы}

\textbf{Mastercard Rules}~--- свод правил, аналогичный
Visa Core Rules. Обновляется дважды в~год. Доступен
на~\domain{mastercard.com}~\sidecite{mc-rules}.

\textbf{Mastercard Chargeback Guide}~--- правила оспаривания
транзакций~\sidecite{mc-chargeback-guide}.

\textbf{Security Rules and Procedures}~--- требования
к~безопасности для~участников сети.

\subsection*{Документы с~ограниченным доступом}

\textbf{Banknet Processing and IPM Technical Manual}~---
форматы авторизационных сообщений (DMSA) и~клиринговых
файлов (IPM). Доступны только для~процессоров.

\textbf{M/Chip Requirements}~--- спецификации
EMV-реализации \scheme{Mastercard}. Доступны
только участникам~\sidecite{mc-mchip}.

\textbf{MDES (Mastercard Digital Enablement Service)}~---
спецификации токенизации~\sidecite{mc-mdes}. Доступны
для~Token Requestors.

\subsection*{Mastercard Developers}

\scheme{Mastercard} Developers
(\domain{developer.mastercard.com})~--- портал
с~открытыми API, sandbox-средами и~документацией.
Здесь~доступны API для~Mastercard Send, Mastercard
Gateway, Open Banking и~других продуктов. Регистрация
бесплатная.

\begin{practice}
  Правила \scheme{Visa} и~\scheme{Mastercard} во~многом
  пересекаются, но~отличаются в~деталях: сроки chargeback,
  reason codes, пороги мониторинга, требования к~3DS~---
  всё~это~специфично для~каждой сети. Не~экстраполируйте
  правила одной сети на~другую: \enquote{у~Visa так~---
    значит, у~Mastercard тоже}~--- частая ошибка, которая
  обнаруживается при~аудите или~при~получении штрафа.
\end{practice}


\section{НСПК и~Мир}
\label{sec:ps-nspk}

Национальная система платёжных карт (НСПК)~---
оператор платёжной системы Мир и~ОПКЦ для~всех
внутрироссийских карточных
транзакций~(см.~главу~\ref{ch:mir}).

\subsection*{Ключевые документы}

\textbf{Правила платёжной системы Мир}~--- основной
нормативный документ, определяющий обязанности участников.
Доступен на~сайте НСПК (\domain{nspk.ru}) в~разделе
документов платёжной системы~\sidecite{nspk-mir-contactless}.

\textbf{Технические спецификации НСПК}~--- форматы сообщений,
требования к~подключению, EMV-конфигурация для~карт Мир,
спецификации МРА. Доступны только для~участников
платёжной системы.

\textbf{Документация по~СБП}~--- правила, API-спецификации,
руководства по~подключению. Основная часть доступна
участникам; базовые описания~--- на~сайте
НСПК~\sidecite{nspk-sbp}.

\subsection*{Особенности}

Документация НСПК менее публична, чем~документация
международных платёжных систем: большая часть технических
спецификаций доступна только участникам по~NDA.
Публичных developer-порталов со~sandbox-средами
на~момент написания книги нет. Для~получения технической
документации необходимо быть участником платёжной системы
или~работать через банк-участник.


\section{Банк России}
\label{sec:ps-cbr}

Регулятор публикует нормативные акты, которые определяют
правовую рамку для~всех участников платёжного
рынка~(см.~главу~\ref{ch:regulation}).

\subsection*{Ключевые документы}

\textbf{\law{161-ФЗ} «О~национальной платёжной системе»}~---
базовый закон, определяющий понятия оператора платёжной
системы, оператора услуг платёжной инфраструктуры,
электронных денежных средств~\sidecite{fz-161}.

\textbf{\law{382-П}}~--- положение о~требованиях к~обеспечению
защиты информации при~осуществлении переводов денежных
средств~\sidecite{cbr-382p}.

\textbf{\law{683-П}}~--- положение об~установлении обязательных
для~кредитных организаций требований к~обеспечению защиты
информации~\sidecite{cbr-683p}.

\textbf{\law{266-П}}~--- положение о~порядке осуществления
наблюдения в~национальной платёжной
системе~\sidecite{cbr-266p}.

\subsection*{Где искать}

Актуальные тексты нормативных актов~--- на~сайте Банка России
(\domain{cbr.ru}) в~разделе \enquote{Нормативные и~иные акты}
и~на~правовых порталах: \domain{pravo.gov.ru} (официальный
текст), \domain{consultant.ru} и~\domain{garant.ru}
(с~комментариями и~историей изменений).

Банк России публикует информационные письма и~разъяснения,
которые не~имеют силы нормативного акта, но~описывают
позицию регулятора по~спорным вопросам. Следить за~ними~---
на~сайте ЦБ в~разделе \enquote{Информационные материалы}.

\begin{practice}
  Нормативные акты ЦБ часто ссылаются друг на~друга,
  образуя сложную сеть зависимостей: \law{382-П} ссылается
  на~\law{161-ФЗ}, \law{683-П} ссылается на~\law{382-П},
  и~так далее. При~чтении любого акта держите под~рукой
  тексты всех связанных документов~--- или~используйте
  правовой портал с~перекрёстными ссылками (Консультант+,
  Гарант), где~можно перейти по~ссылке к~связанному акту
  одним кликом.
\end{practice}


\section{ISO}
\label{sec:ps-iso}

Международная организация по~стандартизации (ISO) публикует
стандарты, широко используемые в~платёжной индустрии,
но~в~отличие от~EMVCo и~PCI~SSC, стандарты ISO~---
платные.

\subsection*{Ключевые стандарты}

\textbf{\iso{8583}}~--- формат сообщений для~финансовых
транзакций, описанный в~главе~\ref{ch:iso8583}. Стандарт
определяет структуру (MTI, bitmap, data elements), но~каждая
платёжная система реализует свой~диалект.

\textbf{\iso{20022}}~--- универсальный стандарт финансовых
сообщений на~базе XML/ASN.1, используемый в~SWIFT MX,
SEPA, БЭСП и~десятках других
систем~(см.~главу~\ref{ch:multicurrency}).

\textbf{ISO~4217}~--- коды валют (643~--- RUB, 978~--- EUR,
840~--- USD); используются в~\de{49}, \de{50}, \de{51}
сообщений \iso{8583}.

\textbf{ISO~3166}~--- коды стран; используются
в~авторизационных сообщениях и~клиринговых файлах.

\textbf{ISO~9564}~--- стандарт управления и~защиты PIN,
определяющий форматы PIN block (ISO~Format~0, 1, 3,
4)~\sidecite{iso-9564}.

\subsection*{Как работать с~платными стандартами}

Стандарты ISO стоят от~100 до~300~швейцарских франков
за~документ, и~для~многих организаций это~не~проблема.
Но~если~доступ ограничен, есть несколько подходов:

\begin{itemize}
  \item Для~\iso{8583}: реальной необходимости покупать
        стандарт обычно нет, потому что~каждая карточная сеть
        публикует свою реализацию (диалект), которая и~является
        рабочим документом для~интеграции. Стандарт полезен
        для~общего понимания, но~не~для~разработки.
  \item Для~\iso{20022}: \domain{iso20022.org} содержит
        полные schema definitions, message definitions
        и~catalogue~--- бесплатно. Это~достаточно для~разработки;
        сам~стандарт нужен только для~формального compliance.
  \item Для~ISO~4217 и~ISO~3166: актуальные списки кодов
        свободно доступны на~сайтах ISO и~многочисленных
        зеркалах.
\end{itemize}


\section{Отслеживание изменений}
\label{sec:ps-tracking}

Знать, где~искать документы,~--- необходимо, но~недостаточно:
документы обновляются, и~изменения могут затронуть ваш~продукт.
Пропущенное обновление правил~--- это~штраф, нарушение
compliance или~отклонение транзакций, которые вчера проходили
нормально.

\subsection*{Бюллетени карточных сетей}

\scheme{Visa} и~\scheme{Mastercard} публикуют бюллетени
(Visa Business News, Mastercard Global Operations Bulletin),
анонсирующие изменения правил, новые продукты, обновления
протоколов и~дедлайны по~миграции. Бюллетени рассылаются
участникам сети; если~вы~не~участник~--- запрашивайте
пересылку через ваш~банк-эквайер или~процессор.

Типичный цикл изменений: бюллетень публикуется за~6--12~месяцев
до~вступления в~силу, давая участникам время на~доработку
систем. Изменения вступают в~силу в~апреле или~октябре
(для~\scheme{Visa}) и~в~аналогичные даты
для~\scheme{Mastercard}.

\subsection*{PCI SSC}

PCI~SSC публикует обновления стандартов с~указанием дат
вступления в~силу (effective date) и~даты окончания
переходного периода (sunset date). Подписка на~PCI
Perspectives Blog и~рассылку PCI~SSC~--- бесплатный
способ отслеживать изменения.

\subsection*{EMVCo}

EMVCo публикует Specification Bulletins~--- точечные изменения
к~спецификациям. Список бюллетеней доступен на~сайте EMVCo;
подписки на~рассылку нет, но~раздел обновляется регулярно.

\subsection*{Банк России}

ЦБ публикует проекты нормативных актов для~общественного
обсуждения на~\domain{regulation.gov.ru} до~их~принятия.
Финальные версии появляются на~\domain{cbr.ru}. Для~отслеживания
полезны правовые системы (Консультант+, Гарант), которые
уведомляют об~изменениях в~отслеживаемых документах.

\subsection*{Отраслевые ассоциации и~сообщества}

Помимо первоисточников, полезны:

\begin{itemize}
  \item \textbf{EAST (European Association for Secure Transactions)}~---
        отчёты о~карточном фроде в~Европе, тренды, статистика.
  \item \textbf{MRC (Merchant Risk Council)}~--- сообщество
        специалистов по~fraud prevention, вебинары, исследования.
  \item \textbf{Payments Cards \& Mobile (PCM)}~--- отраслевое
        издание, обзоры регуляторных изменений.
  \item \textbf{PLUSWORLD}~--- русскоязычное отраслевое издание,
        покрывающее российский и~СНГ платёжный рынок.
\end{itemize}

\begin{important}
  Отраслевые издания и~вебинары~--- хорошие сигнальные
  источники: они~помогают узнать \emph{о~чём} изменение
  и~когда оно~вступает в~силу. Но~для~принятия решений~---
  доработки системы, изменения flow, подготовки к~аудиту~---
  всегда обращайтесь к~первоисточнику: бюллетеню сети,
  тексту стандарта, нормативному акту. Пересказ~--- даже
  качественный~--- может упустить нюанс, который окажется
  критичным для~вашего~случая.
\end{important}


\section*{Итоги главы}

\begin{itemize}
  \item EMVCo публикует спецификации бесплатно (Book~1--4,
        Contactless, Tokenisation, 3DS); начинайте чтение
        с~Book~4 (пользовательские сценарии), затем~Book~3
        (Transaction Flow и~Data Elements).
  \item PCI~DSS, SAQ, PIN Security, P2PE доступны бесплатно
        на~сайте PCI~SSC; при~чтении PCI~DSS начинайте
        с~определения scope и~используйте колонку Guidance.
  \item Правила \scheme{Visa} и~\scheme{Mastercard} публичны
        (Core Rules / Mastercard Rules), но~технические
        спецификации (TADG, BASE~I/II, IPM) доступны
        только участникам; developer-порталы обеих сетей~---
        бесплатный источник API-документации и~sandbox.
  \item У~НСПК большая часть техдокументации закрыта.
  \item Нормативные акты ЦБ~РФ ищите на~\domain{cbr.ru},
        \domain{pravo.gov.ru} и~правовых порталах.
  \item Стандарты ISO (\iso{8583}, \iso{20022}) платные,
        но~рабочие спецификации~--- диалекты сетей и~schema
        definitions \iso{20022}~--- доступны бесплатно.
  \item Отслеживайте изменения через бюллетени сетей,
        PCI~SSC Blog, EMVCo Bulletins и~правовые системы;
        для~решений~--- всегда обращайтесь к~первоисточнику,
        а~не~к~пересказу.
\end{itemize}
